\section{Étape 3 -- Configuration réseau du coupleur}
\label{secEtapeConfigReseau}

Cette étape consiste à attribuer au coupleur les paramètres réseau
(adresse IP, masque, passerelle) qui lui permettront de communiquer avec
les équipements qu'il doit scruter.

\subsection{Accès à la configuration réseau}

Ouvrir le navigateur de DTM (menu \emph{Outils -> Navigateur de DTM}), puis
double-cliquer sur le réseau concerné pour faire apparaître sa fenêtre de
configuration.

\subsection{Paramétrage de la voie}

Dans la rubrique \textbf{TCP/IP} de cette fenêtre, renseigner :

\begin{itemize}
    \item l'\textbf{adresse IP} du coupleur, conforme au plan d'adressage
    de l'installation ;
    \item le \textbf{masque de sous-réseau} ;
    \item la \textbf{passerelle}, le cas échéant (une valeur
    \texttt{0.0.0.0} indique l'absence de passerelle).
\end{itemize}

\UPSTIremarque[Exemple -- réseau NOC\_ETHIP\_CONVOYEUR]{
Adresse IP : \texttt{172.16.180.180} ; masque : \texttt{255.255.255.0} ;
passerelle : \texttt{172.16.180.254}.
}

\UPSTIremarque[Exemple -- réseau NOC\_ETHIP\_ROBOT]{
Adresse IP : \texttt{192.168.0.1} ; masque : \texttt{255.255.255.0} ; pas de
passerelle (\texttt{0.0.0.0}). Sur cette installation, la configuration par
défaut du coupleur répond déjà à cette attente : il convient néanmoins de
la vérifier systématiquement.
}
